%% ============================================================
%% 分片 19：原文第 7 节「Fluxes and $S$-Duality」（原文 glshort.tex 起于第 5685 行）
%% 由 tools/harvmac2tex.py 从 arXiv 源码（hep-th/0604151v3）自动生成。
%% 只译英文散文；公式、编号、\label、\cite 一律不改。
%% ============================================================

\section{通量与 $S$-对偶}

\label{sec:centertop}

\begin{keytakeaways}
  \begin{itemize}[leftmargin=1.4em,itemsep=2pt,topsep=2pt]
    \item \textbf{物理图像}：规范群的中心 $\CZ=\pi_1(G_{\rm ad})$ 是有限阿贝尔群，并且具有自然的自对偶；这使离散磁通量 ${\bf m}$ 与离散电通量 ${\bf e}$ 可以取遍互为对偶的有限群 ${\bf M}$、${\bf E}$，从而让 $S$-对偶在量子理论中得以实现。
    \item \textbf{数学工具}：把理论放在 $M=S^1\times W$ 上，$\Gad$-丛的特征类 $\xi(E)$ 分解为 ${\bf a}\in H^1(W,\CZ)$ 与 ${\bf m}\in H^2(W,\CZ)$ 两部分；庞加莱对偶给出完美配对 $H^1(W,\CZ)\times H^2(W,\CZ)\to U(1)$，故二者互为对偶的有限阿贝尔群。
    \item \textbf{关键结论}：固定 ${\bf m}$ 有 Hilbert 空间解释，固定 ${\bf a}$ 则没有，因此须对 ${\bf a}$ 作傅里叶变换、引入特征标 ${\bf e}$，得到 ${\cal H}_{{\bf e},{\bf m}}$；$S$-对偶交换 ${\bf E}$ 与 ${\bf M}$，把 $({\bf e},{\bf m})$ 送到 $(-{\bf m},{\bf e})$，并把 $\tau$ 送到 $-1/\tau$。
    \item \textbf{边界与展望}：泛丛对伴随群 $G_{\rm ad}$ 存在，对单连通覆盖 $\bar G$ 则被非平凡类 $\zeta=\xi(\EUE_{\rm ad})\in H^2(\M,\CZ)$ 所阻碍；但它总可以作为扭化向量丛存在。
  \end{itemize}
\end{keytakeaways}

本节我们将考察与规范群中心有关的现象，这些现象对于理解 Hitchin 模空间的 $S$-对偶以及几何 Langlands 纲领是相关的。

我们先在第~\ref{sec:reviewthis}~节回顾离散电通量与磁通量以及 $S$-对偶对它们的作用。更多细节可参看 \cite{ref45} 第 3 节以及 \cite{ref26}。关于规范场论中的离散电、磁通量，原始文献是 \cite{ref101}。我们还会解释一些数学概念，例如泛丛的概念。然后在第~\ref{sec:comptwop}~节展开我们的应用。

\subsection{回顾}

\label{sec:reviewthis}

本节中，$G$ 表示紧半单李群；$\bar G$ 是它的万有覆盖；$G_{\rm {ad}}$ 是对应的伴随群。$\bar G$ 的中心（它同时也是 $G_{\rm {ad}}$ 的基本群）是一个有限阿贝尔群 $\CZ$。对于单的 $G$，除 $\bar G=Spin(4n)$ 外，$\CZ$ 在所有情形下都是循环群，如表中所列。

从表中可以注意到一个令人欣慰的事实：对于非单连通型（non-simply-laced）的群，$S$-对偶的某些微妙之处无关紧要，因为 $G_2$ 与 $F_4$ 的中心是平凡的。

\begin{yizhu}[术语]
按：simply-laced（各根等长，A、D、E 型）本文译作「单连通型」，与 simply connected（单连通的）只差一字；两者英文原词不同，阅读时须按上下文区分。本节的「单连通覆盖 $\bar G$」指的是后者。
\end{yizhu}

\par\medskip

\begin{table}[htbp]
\centering
\begin{tabular}{ll}
\toprule
群 & 中心 \\
\midrule
$SU(N)$ & $\Bbb{Z}_N$ \\
$Spin(2n+1)$ & $\Bbb{Z}_2$ \\
$Spin(4n+2)$ & $\Bbb{Z}_4$ \\
$Spin(4n)$ & $\Bbb{Z}_2\times\Bbb{Z}_2$ \\
$Sp(n)$ & $\Bbb{Z}_2$ \\
$E_6$ & $\Bbb{Z}_3$ \\
$E_7$ & $\Bbb{Z}_2$ \\
\bottomrule
\end{tabular}
\caption{具有非平凡中心的单连通单李群。表中略去的 $F_4$、$G_2$ 与 $E_8$ 有平凡中心。（原文 Table~3；原文此处误作 $F_2$、$G_4$）}
\label{tab:centers}
\end{table}

一个重要性质是：群 ${\cal Z}$ 自然地自对偶。阿贝尔群 $B$ 的对偶是 $B^\vee={\rm Hom}(B,U(1))$。有限阿贝尔群总与它自己的对偶同构，但并不是自然地同构。然而 $\bar G$ 的中心 $\CZ$ 总具有自然的自对偶。这一对偶最好用同态 % 式 (7.1)
\setcounter{equation}{0}
\begin{equation}\label{eq:plook}
\Upsilon:\CZ\times \CZ\to
U(1),
\end{equation}
来表达，它是对称的，即 $\Upsilon(a,b)=\Upsilon(b,a)$，并且使得 $\CZ$ 到 $U(1)$ 的任一同态都可写成 $a\to \Upsilon(a,b)$，其中 $b$ 唯一。这样的 $\Upsilon$ 称为「完美配对」。$\Upsilon$ 的构造如下。若 $G$ 是单连通型，则 $\CZ=\Lambda^\vee/\Lambda$，其中 $\Lambda$ 是 $\bar G$ 的根格，$\Lambda^\vee$ 是权格。此时配对 $\Upsilon$ 定义为
% 式 (7.2)
\setcounter{equation}{1}
\begin{equation}\label{eq:rofo}
\Upsilon(a,b)=\exp(2\pi i\langle a,b\rangle),
\end{equation}
其中 $\langle~,~\rangle$ 是 $\Lambda^\vee$ 上通常的二次型（它在 $\Lambda$ 上的限制是偶的、整值的，并在 $\Lambda^\vee$ 上取有理值）。若 $G$ 不是单连通型，则它的中心平凡或为 $\Bbb{Z}_2$，于是恰好容许一种自对偶，我们称之为 $\Upsilon$。

$M$ 上的 $\Gad$-丛 $E$ 有一个特征类 $\xi(E)$，取值于 $H^2(M,\CZ)$。例如若 $G_{{\rm ad}}=SO(3)$，则 $\xi(E)$ 是第二 Stieffel-Whitney 类 $w_2(E)$。对每一种 $\xi$ 的选取，都可以通过把路径积分限制到具有给定 $\xi$ 的丛上来定义配分函数 $Z_\xi$。这些 $Z_\xi$ 按 $S$-对偶群的幺正表示变换 \cite{ref45,ref26}，但这一点我们在此不作解释。

\begin{yizhu}[原文疑误]
按：原文拼作 Stieffel-Whitney，通行拼法为 Stiefel-Whitney（施蒂费尔-惠特尼类）；这里指的是第二 Stiefel-Whitney 类 $w_2(E)$。
\end{yizhu}

我们转而集中于哈密顿描述，并取 $M=S^1\times W$，其中 $W$ 是三维流形。此时有 $H^2(M,\CZ)=H^1(W,\CZ)\oplus H^2(W,\CZ)$，因此 $\xi(E)$ 可以分解为 $\xi={\bf a}\oplus {\bf m}$，其中 ${\bf a}\in H^1(W,\CZ)$，${\bf m}\in H^2(W,\CZ)$。

\begin{yizhu}[物理图像]
按：这里的「通量」是拓扑（离散）数据，而不是场强通量：${\bf m}\in H^2(W,\CZ)$ 是绕 $W$ 中闭曲面的磁通量，${\bf a}\in H^1(W,\CZ)$ 是沿 $S^1$ 方向、与 $W$ 中闭曲线相配对的对应量；二者是同一个特征类 $\xi(E)$ 的两个分量。参见 Witten 2008 综述中关于离散通量与对偶群的讨论。
\end{yizhu}

考虑配对 % 式 (7.3)
\setcounter{equation}{2}
\begin{equation}\label{eq:junko}
\hat \Upsilon:H^2(M,\CZ)\times
H^2(M,\CZ)\to H^4(M,U(1))=U(1)
\end{equation}
它由把上积 $H^2(M,\CZ)\times H^2(M,\CZ)\to H^4(M,\CZ\times\CZ)$ 与映射 $\Upsilon:\CZ\times\CZ\to U(1)$ 复合而得。按庞加莱对偶，$\hat\Upsilon$ 是完美配对，从而使 $H^2(M,\CZ)$ 成为自对偶阿贝尔群。更一般地，对 $n$ 维闭定向流形 $Y$，庞加莱对偶连同 $\CZ$ 的自对偶给出完美配对
% 式 (7.4)
\setcounter{equation}{3}
\begin{equation}\label{eq:unko}
H^d(Y,\CZ)\times H^{n-d}(Y,\CZ)\to U(1), ~~0\leq d\leq
n.
\end{equation}
特别地，我们有完美配对
% 式 (7.5)
\setcounter{equation}{4}
\begin{equation}\label{eq:normo}
\bar\Upsilon:H^1(W,\CZ)\times H^2(W,\CZ)\to U(1).
\end{equation}
所以它们是互为对偶的阿贝尔群，特别地
% 式 (7.6)
\setcounter{equation}{5}
\begin{equation}\label{eq:junok}
H^2(W,\CZ)={\rm Hom}(H^1(W,\CZ),U(1)).
\end{equation}

现在我们来回顾 $S$-对偶如何在哈密顿框架中实现。指定了 ${\bf m}$ 的规范场论配分函数有自然的 Hilbert 空间解释，因为 ${\bf m}$ 是初始时刻数据的一部分，并且与时间无关。然而，由于 ${\bf a}$ 不能由固定时刻的数据表达出来，固定 ${\bf a}$ 的配分函数不能赋予哈密顿解释。于是必须引入有限阿贝尔群 $H^1(M,\CZ)$ 的一个特征标 ${\bf e}$，即同态
% 式 (7.7)
\setcounter{equation}{6}
\begin{equation}\label{eq:kilno}
{\bf e}:H^1(W,\CZ)\to U(1).
\end{equation}
对每一组 ${\bf e}$ 与 ${\bf m}$ 的选取，和 $\sum_{{\bf a}}{\bf e}({\bf a})Z_{{\bf a},{\bf m}}$ 都可以解释为某个 Hilbert 空间 ${\cal H}_{{\bf e},{\bf m}}$ 中的迹。这个和有 Hilbert 空间解释，是因为以权因子 ${\bf e}({\bf a})$ 对 ${\bf a}$ 求和与沿 $S^1$ 方向的切割、粘合（量子场论中通常称为集团分解）相容。${\bf e}$ 与 ${\bf m}$ 分别称为离散电通量与离散磁通量 \cite{ref101}。由 \eqref{eq:junok}，${\bf e}$ 也可以等价地看作 $H^2(W,\CZ)$ 中的一个元素，正如 ${\bf m}$ 一样。

\begin{yizhu}[原文疑误]
按：原文正文说特征标 ${\bf e}$ 是 $H^1(M,\CZ)$ 的同态，而紧随其后的式 (7.7) 写作 $H^1(W,\CZ)\to U(1)$。既然 ${\bf a}\in H^1(W,\CZ)$，对偶的特征标应定义在 $H^1(W,\CZ)$ 上；原文此处疑把 $W$ 写成了 $M$。译文照原文保留了这一措辞。
\end{yizhu}

于是对每一组
% [WARN] 公式 \eqn\eqref{eq:ytro} 无编号或重复
\begin{equation*}
\tag*{(7.8)}
\begin{aligned}{\bf e}&\in {\bf E}={\rm
Hom}(H^1(W,\CZ), U(1))\\
                             {\bf m}& \in {\bf M}=H^2(W,\CZ).
\end{aligned}
\end{equation*}
的选取，我们都得到一个 Hilbert 空间 ${\cal H}_{{\bf m},{\bf e}}$。所幸由 \eqref{eq:junok}，有限群 ${\bf E}$ 与 ${\bf M}$ 同构。这使得 $S$-对偶对称性成为可能。对偶变换 $S$ 交换 ${\bf E}$ 与 ${\bf M}$，并把 $({\bf e},{\bf m})\to (-{\bf m},{\bf e})$：
% 式 (7.9)
\setcounter{equation}{8}
\begin{equation}\label{eq:coco}
S:{\cal H}_{{\bf e},{\bf m}}\to {\cal H}_{-{\bf m},{\bf
e}}.
\end{equation}
当然，它还把 $\tau\to -1/\tau$。（式 \eqref{eq:coco} 是 \eqref{eq:xilj} 的离散类比，其中 $(\vec{\eurm m},\vec{\eurm n})$ 换成了它们的离散对应物 $({\bf m},{\bf e})$。）

\begin{yizhu}[物理图像]
按：$S$ 把固定 $({\bf e},{\bf m})$ 的扇区送到 $(-{\bf m},{\bf e})$，即交换电、磁通量并带一个符号；这与连续理论中 $S:\tau\to-1/\tau$ 以及 $(E,B)\to(B,-E)$ 的变换互为影子。参见 Witten 2008 综述。
\end{yizhu}

虽然可以定义所有这些 Hilbert 空间 ${\cal H}_{{\bf e},{\bf m}}$，但给定的规范场论构造未必用到全部。若我们在 $M=S^1\times W$ 上以单连通规范群 $\bar G$ 作规范场论，则必须把进入配分函数 $Z_{{\bf a},{\bf m}}(M)$ 的特征类 ${\bf a}$ 与 ${\bf m}$ 置零。在 Hilbert 空间层面，这意味着 ${\bf e}$（它由关于 ${\bf a}$ 的傅里叶变换产生）任取，而 ${\bf m}=0$。于是 % 式 (7.10)
\setcounter{equation}{9}
\begin{equation}\label{eq:ytoto}
{\cal H}(\bar G,W)=\oplus_{{\bf e}}{\cal
H}_{{\bf e},0}.
\end{equation}
计及假定下的 ${\bf e},{\bf m}$ 在 $S$-对偶下的变换律，我们看到对偶理论的 Hilbert 空间必为 % 式 (7.11)
\setcounter{equation}{10}
\begin{equation}\label{eq:momjo}
\oplus_{{\bf
m}}{\cal H}_{0,{\bf m}}.
\end{equation}
这个 Hilbert 空间与 $M=S^1\times W$ 上的路径积分有何关系？令 ${\bf e}$ 为零意味着对 ${\bf a}$ 求和。因此在变量 ${\bf a},{\bf m}$ 的表述下，\eqref{eq:momjo} 意味着必须以相同的权重对所有 ${\bf a},{\bf m}$ 求和。换言之，在对偶理论中我们对具有一切可能 $\xi$ 的 $^L\Gad$-丛求和。于是对偶理论的规范群为 $^L\Gad$，正如所期。

\par\medskip\emph{泛丛}

我们还需要另一项背景知识，即泛丛的概念。

先考虑黎曼面 $C$ 上 $G$-丛的模空间 $\M(G,C)$。对每个点 $p\in \M(G,C)$，都有相应的一个由 $p$ 参数化的 $G$-丛 $E_p\to C$；它在同构意义下由 $p$ 确定。泛丛（若存在）是一个 $G$-丛 ${\EUE} \to \M\times C$，使得对任意 $p\in \M$，${\EUE}$ 限制到 $p\times C$ 上都同构于 $E_p$。规范场论给出 \cite{ref91} 一个分析泛丛的非常直接的途径，但我们在此不作解释。

朴素地，我们只需这样定义 ${\EUE}$：使其限制到 $p\times C$ 上就是由 $p$ 决定的 $C$ 上「那个」丛。问题之所以出现，是因为首先 $p$ 只决定 $C$ 上 $G$-丛的同构类，而不是一个具体的丛。这在局部不会造成困难\footnote{这一陈述在 $\M$ 的光滑点附近成立。比本文更精细的分析表明，阻碍泛丛的类 $\zeta(\EUE_{\rm ad})$ 可以在奇点处有局部贡献。}；我们只需在某个特定的 $p$ 处取同构类中正确的一个丛，然后在小邻域内变形它。于是对每个小开集 $U_i\subset \M$，我们都可以选取一个泛丛 ${\EUE}_i\to U_i\times C$，而且 ${\EUE}_i$ 在同构意义下唯一。

下一步是把这些 ${\EUE}_i$ 在 $(U_i\cap U_j)\times C$ 上粘起来。${\EUE}_i$ 与 ${\EUE}_j$ 在 $(U_i\cap U_j)\times C$ 上同构，所以我们在这个交叠上选取同构 $\Theta_{ij}:{\EUE}_i\cong {\EUE}_j$，并令 $\Theta_{ji}=\Theta_{ij}^{-1}$。

现在，如果在三重交叠 $(U_i\cap U_j\cap U_k)\times C$ 上复合 $\Theta_{ki}\Theta_{jk}\Theta_{ij}$ 等于 1，那么我们就可以相容地把这些 ${\EUE}_i$ 粘起来，得到所要的泛丛 ${\EUE}$。

若规范群是伴随群 $G_{{\rm ad}}$，粘合没有问题。一般的 $G_{{\rm ad}}$-丛没有自同构（例外出现在 $\M$ 的奇点处）。因此 $\Theta_{ki}\Theta_{jk}\Theta_{ij}$ 由于是 ${\EUE}_i$（限制到 $(U_i\cap U_j\cap U_k)\times C$）到自身的同构，在三重交叠的一般点上等于恒等，从而处处等于恒等。

于是，泛 $G_{{\rm ad}}$-丛 ${\EUE}_{\rm {ad}}$ 确实存在。现在我们把规范群换成单连通覆盖 $\bar G$（更一般地，换成 $G_{\rm{ad}}$ 的任意具有非平凡中心的覆盖）来重复这一讨论。我们从同样的方式开始。对每个小开集 $U_i\in \M$，选取 $U_i\times C$ 上的泛丛 $\bar {\EUE}_i$，并试图把它们粘成泛 $\bar G$-丛 $\bar {\EUE}$。

结果之所以不同，是因为一般的 $\bar G$-丛确实有非平凡的自同构群；中心 $\CZ$ 是任意 $\bar G$-丛 $E$ 的一个自同构群。

因此在上面的论证中，三重交叠上的复合映射 $\Theta_{ki}\Theta_{jk}\Theta_{ij}$ 未必是恒等；而是 % 式 (7.12)
\setcounter{equation}{11}
\begin{equation}\label{eq:yrf}
\Theta_{ki}\Theta_{jk}\Theta_{ij}=f_{ijk},
\end{equation}
其中 $f_{ijk}\in\CZ$。这些 $f_{ijk}$ 组成一个二维闭链，定义了一个元素 $\zeta\in H^2(\M,\CZ)$。已知 $\zeta$ 是非平凡的；它是泛 $\bar G$-丛 $\bar{\EUE}$ 存在性的障碍。

\begin{yizhu}[易错点]
按：障碍 $\zeta$ 是 $H^2(\M,\CZ)$ 中的类，而不是 $H^1$ 中的类；它来自三重交叠上 $\Theta_{ki}\Theta_{jk}\Theta_{ij}\in\CZ$ 的粘合失败。$G_{\rm ad}$-丛没有自同构，故无此问题；$\bar G$-丛的全纯中心 $\CZ$ 自同构正是根源。
\end{yizhu}

还有另一种方式解释 $\zeta$ 的存在。设 $\M(G_{\rm{ad}},C)_0$ 是 $\M(G_{\rm{ad}},C)$ 中参数化拓扑平凡丛的分支——这些丛限制到 $p\times C$（$p\in \M(G_{\rm{ad}},C)$）上时可以提升为 $\bar G$-丛。考虑 $\M(G_{\rm{ad}},C)_0\times C$ 上的泛丛 ${\EUE}_{\rm{ad}}$。如果我们能把 ${\EUE}$ 提升为 $\bar G$-丛 $\bar {\EUE}$，那么（把它拉回到 $\M(\bar G,C)\times C$，后者是 $\M(G_{\rm {ad}},C)_0\times C$ 的有限覆盖）这就是泛 $\bar G$-丛。一般地，把 $G_{\rm ad}$-丛 ${\EUE}_{\rm{ad}}$ 提升为 $\bar G$-丛 $\bar {\EUE}$ 的障碍是特征类 $\xi({\EUE}_{\rm{ad}})$。所以泛 $\bar G$-丛存在性的障碍是 $\zeta=\xi({\EUE}_{\rm {ad}})$。

对 Higgs 丛有类似的故事。泛 Higgs 丛是一个对 $( {\EUE},\hat \varphi)$，其中 $ {\EUE} $ 是 $\MH(G,C)\times C$ 上的 $G$-丛，$\hat\varphi\in H^0( \MH(G,C)\times C,{\rm {ad}}({\EUE})\otimes K_C)$，并满足如下条件：对每个 $p\in \MH(G,C)$，$( {\EUE},\hat\varphi)$ 限制到 $p\times C$ 上应当同构于由 $p$ 参数化的 Hitchin 对 $(E_p,\varphi_p)$。

把已经给出的论证原封不动地搬过来，即可证明对规范群 $G_{\rm{ad}}$，泛 Higgs 丛 $( {\EUE}_{\rm {ad}},\hat\varphi)$ 确实存在。但对单连通规范群 $\bar G$（或 $G_{\rm {ad}}$ 的任一非平凡覆盖），泛 Higgs 丛不存在。其障碍在于：一般的 Hitchin 对 $(E,\varphi)$（其中 $E$ 为 $\bar G$-丛）具有自同构群 $\CZ$。这个障碍正是把泛 $G_{{\rm ad}}$ Higgs 丛 $({\EUE}_{\rm{ad}},\hat\varphi)$ 的「丛」部分提升为 $\bar G$-丛的障碍。因此它是
% 式 (7.13)
\setcounter{equation}{12}
\begin{equation}\label{eq:plooo}
\zeta=\xi({\EUE}_{\rm{ad}}).
\end{equation}

\par\medskip\emph{泛丛作为扭化向量丛}

虽然对非伴随型的 $G$，泛 $G$-丛作为向量丛或主丛不存在，但它作为扭化向量丛确实存在。我们将以一种非常初等的方式来描述这个概念。

构造流形 $X$ 上秩为 $N$ 的普通向量丛 $V$ 的一种办法是：用小的开集 $U_i$ 覆盖 $X$，在每个开集上选取秩 $N$ 的平凡丛 $V_i\cong U_i\times \Bbb{C}^N$。然后在交叠 $U_i\cap U_j$ 上用粘合映射 $v_{ij}:V_i\cap V_j\to U(N)$ 把 $V_i$ 粘到 $V_j$，并令 $v_{ji}=v_{ij}^{-1}$。如果在三重交叠上有 % 式 (7.14)
\setcounter{equation}{13}
\begin{equation}\label{eq:nofto}
v_{ki}v_{jk}v_{ij}=1,
\end{equation}
那么这些 $V_i$ 就可以相容地粘成一个秩 $N$ 的向量丛 $V\to X$。

现在假设给定 $b\in H^2(X,U(1))$。则 $b$ 可以由三重交叠 $U_i\cap U_j\cap U_k$ 上定义的 $U(1)$-值闭链 $b_{ijk}$ 显式表示。扭化向量丛由同样类型的数据 $v_{ij}:V_i\cap V_j\to U(N)$ 定义。但现在我们要求的不再是 \eqref{eq:nofto}，而是
% 式 (7.15)
\setcounter{equation}{14}
\begin{equation}\label{eq:hofto}
v_{ki}v_{jk}v_{ij}=b_{ijk}.
\end{equation}
于是，$b$-扭化向量丛 $V$ 不是通常意义上的向量丛。然而，伴随向量丛 ${\rm ad}(V)$ 是普通向量丛，因为一旦转到伴随表示，\eqref{eq:hofto} 中的相位 $b_{ijk}$ 就消失了。

现在回到构造泛丛的问题。本文一般对下面一点稍欠精确：所谓 $G$-丛，究竟是指纤维为 $G$ 的副本的主 $G$-丛，还是在 $G$ 的某个忠实表示下的相伴向量丛。主丛使得对任意 $G$ 都适用的统一分析成为可能；但对像 $U(N)$ 这样具有方便的忠实表示（$N$ 维表示）的群，用向量丛来思考是有益的。

在讨论泛丛时，更精确一些是有帮助的。我们把上面讨论中的过渡函数 $\Theta_{ij}$ 解释为 $\bar G$-值函数，不选取任何特定的表示。它们到 $G_{ad}$ 的投影给出泛主 $G_{ad}$-丛 ${\EUE}_{\rm ad}$ 的过渡函数，但由于关系
% 式 (7.16)
\setcounter{equation}{15}
\begin{equation}\label{eq:flutter}
\Theta_{ki}\Theta_{jk}\Theta_{ij}=f_{ijk},
\end{equation}
其中 $f_{ijk}$ 定义了类 $\zeta=\xi({\EUE}_{ad})\in H^2(\MH(G_{\rm ad},C),\CZ)$，我们无法把它提升为主 $\bar G$-丛。现在取 $\bar G$ 的一个不可约表示 $\varrho$。我们把 $\Theta_{ij}$ 在表示 $\varrho$ 下求值所得的过渡函数记作 $\Theta_{ij}^\varrho$。同样把 $f$ 在表示 $\varrho$ 下求值记作 $\varrho(f)$。在不可约表示 $\varrho$ 中，$\bar G$ 的中心按纯量乘法作用，因此 $\varrho(f)$ 取值于 $U(1)$。我们有
% 式 (7.17)
\setcounter{equation}{16}
\begin{equation}\label{eq:glutter}
\Theta^\varrho_{ki}\Theta^\varrho_{jk}\Theta^\varrho_{ij}=\varrho(f_{ijk}).
\end{equation}

这些量 $\varrho(f_{ijk})$ 是一个闭链，定义了元素 $\varrho(\zeta)\in H^2(\MH(G_{\rm ad},C),U(1))$。

若 $\varrho(f)=1$，则这些对象 $\Theta^\varrho$ 就是过渡函数，它们定义了一个向量丛 ${\EUE}_\varrho\to \MH(\bar G,C)\times C$，我们可以称之为表示 $\varrho$ 中的泛丛。一般情形下并非如此。然而，比较式 \eqref{eq:flutter} 与 \eqref{eq:hofto} 可以看出，虽然 ${\EUE}_\varrho$ 可能不作为普通向量丛存在，但它总是作为扭化向量丛存在，其扭化为 $\varrho(\zeta)\in H^2(\MH(G_{\rm ad},C),U(1))$。

\begin{yizhu}[易错点]
按：扭化向量丛的粘合只在三重交叠上差一个 $U(1)$ 相位 $b_{ijk}$，因此它不是通常意义上的丛；但取伴随丛时这个相位消失，故 ${\rm ad}(V)$ 仍是普通向量丛。对泛丛而言，取不可约表示 $\varrho$ 后相位变为 $\varrho(f_{ijk})$，于是 ${\EUE}_\varrho$ 是以 $\varrho(\zeta)$ 扭化的向量丛，而不是普通丛。
\end{yizhu}
